\chapter{Enlarged Spleen (Splenomegaly)}

\section*{Definition}
Splenomegaly is an enlargement of the spleen beyond its normal size (approximately 12 cm in length, 150-200 g in weight). It is a sign of underlying disease rather than a primary condition.

\section*{What Happens in Your Body}
The spleen serves as a filter for blood, removing old or abnormal red blood cells and pathogens. Splenomegaly occurs through: (1) increased workload (hemolytic anemias, infections); (2) congestion from portal hypertension; (3) infiltration by abnormal cells (leukemia, lymphoma); (4) extramedullary hematopoiesis.

\section*{Causes}
\begin{itemize}
\item Portal hypertension (cirrhosis, portal vein thrombosis)
\item Infections (mononucleosis, malaria, endocarditis, tuberculosis, HIV)
\item Hematologic disorders (hemolytic anemias, sickle cell disease)
\item Myeloproliferative neoplasms (myelofibrosis, CML)
\item Lymphoproliferative disorders (lymphoma, CLL)
\item Infiltrative diseases (sarcoidosis, amyloidosis, Gaucher disease)
\item Autoimmune diseases (lupus, rheumatoid arthritis)
\item Splenic cysts or abscesses
\item Splenic infarction
\end{itemize}

\section*{When to See a Doctor}
\begin{itemize}
\item Left upper quadrant abdominal pain or fullness
\item Early satiety (spleen compressing the stomach)
\item Palpable spleen on abdominal examination
\item Unexplained weight loss, fever, or night sweats
\item Fatigue, pallor, or easy bruising
\item Signs of portal hypertension (ascites, jaundice)
\item Signs of infection (fever, lymphadenopathy)
\end{itemize}

\section*{Related Symptoms}
\begin{itemize}
\item Left upper quadrant pain
\item Abdominal fullness or bloating
\item Early satiety
\item Fatigue
\item Jaundice
\item Lymphadenopathy
\item Ascites
\item Fever
\end{itemize}

\section*{Self-Care/Home Management}
\begin{itemize}
\item Avoid contact sports or abdominal trauma
\item Avoid heavy lifting or straining
\item Use a seatbelt to protect the abdomen
\item Take acetaminophen for mild discomfort (avoid NSAIDs)
\item Monitor for signs of splenic rupture
\end{itemize}

\section*{How Your Doctor Will Evaluate You}
\begin{itemize}
\item Abdominal ultrasound to confirm splenomegaly
\item Complete blood count with differential
\item Liver function tests
\item Serologic tests for infectious causes
\item CT or MRI of the abdomen
\item Bone marrow biopsy if hematologic malignancy suspected
\end{itemize}

\section*{Treatment Options}
\begin{itemize}
\item Treat the underlying cause
\item Splenectomy for massive splenomegaly causing symptoms
\item Splenic embolization (less invasive alternative)
\item Vaccination before splenectomy
\item Avoid contact sports to prevent splenic rupture
\item Folic acid supplementation in hemolytic anemias
\end{itemize}
